\unitchapter{2}{6}{Software Engineering}{p2-06}

\unitmeta{Expected questions: 8--10 · Target: 9+ · Type: mostly theory (Pressman)
+ metrics numericals (COCOMO, FP, cyclomatic) --- most scoring unit}

\begin{sylbox}

\begin{itemize}
\tightlist
\item[$\square$]
  Software process models: waterfall, prototype, evolutionary (incremental, spiral), concurrent, component-based, formal methods, AOSD, unified process, Agile (XP, Scrum, DSDM, Crystal, FDD, ASD, Kanban)
\item[$\square$]
  Software requirements: functional \& non-functional, eliciting, analysis, requirements specification, SRS, validation, management
\item[$\square$]
  Software design: abstraction, architecture, patterns, modularity, cohesion \& coupling, information hiding, functional independence, refinement, architectural \& UI design
\item[$\square$]
  Software quality: McCall, ISO 9126, quality control \& assurance, risk management, RMMM, software reliability
\item[$\square$]
  Estimation \& scheduling: size (LOC, FP), COCOMO, project scheduling \& staffing, timeline charts
\item[$\square$]
  Software testing: verification \& validation, error/fault/defect/failure, unit/integration/system/acceptance, white-box \& black-box, regression, performance, stress, alpha/beta
\item[$\square$]
  Configuration management; maintenance \& re-engineering; reverse engineering
\end{itemize}

\end{sylbox}

\needspace{131pt}

\hypertarget{p2-06--1-software-process-models}{%
\section{Software Process Models}\label{p2-06--1-software-process-models}}

\needspace{95pt}

\hypertarget{p2-06--11-waterfall-royce-1970}{%
\subsection{Waterfall (Royce, 1970)}\label{p2-06--11-waterfall-royce-1970}}

\diagram{figures/p2-06-00}

\begin{itemize}
\tightlist
\item
  Linear sequential; each phase completes before next. Good when \textbf{requirements are clear \& stable}. No working software until late; high risk.
\item
  \textbf{V-Model}: verification phases (left) paired with validation/testing phases (right): Requirements↔Acceptance test, System design↔System test, Architecture↔Integration test, Module design↔Unit test.
\end{itemize}

\begin{listingblock}{86pt}
\begin{Verbatim}[fontsize=\fontsize{9.0}{10.9}\selectfont]
 Requirements ─────────────────────────► Acceptance testing
   System design ───────────────────► System testing
     Architecture design ────────► Integration testing
        Module design ─────────► Unit testing
                    ╲         ╱
                      Coding
\end{Verbatim}
\end{listingblock}

\needspace{161pt}

\hypertarget{p2-06--12-prototype-model}{%
\subsection{Prototype Model}\label{p2-06--12-prototype-model}}

Quick design → build prototype → customer evaluation → refine → (throwaway or evolutionary). Best when \textbf{requirements are unclear}.

\needspace{95pt}

\hypertarget{p2-06--13-incremental--rad}{%
\subsection{Incremental \& RAD}\label{p2-06--13-incremental--rad}}

\begin{itemize}
\tightlist
\item
  \textbf{Incremental}: deliver in increments; first increment = core product.
\item
  \textbf{RAD} (Rapid Application Development): 60--90 days, component reuse, multiple teams; needs modular \& well-understood requirements.
\end{itemize}

\needspace{220pt}

\hypertarget{p2-06--14-spiral-model-barry-boehm-1986--risk-driven}{%
\subsection{\texorpdfstring{Spiral Model (Barry Boehm, 1986) --- \textbf{risk-driven}}{Spiral Model (Barry Boehm, 1986) --- risk-driven}}\label{p2-06--14-spiral-model-barry-boehm-1986--risk-driven}}

\begin{listingblock}{170pt}
\begin{Verbatim}[fontsize=\fontsize{8.2}{10.0}\selectfont]
            1. Determine objectives,      │      2. Evaluate alternatives,
               alternatives, constraints  │         IDENTIFY & RESOLVE RISKS
                                          │
             ┌────────────────────────────┼────────────────────────────┐
             │       ┌────────────────────┼────────────────────┐       │
             │       │       ┌────────────┼────────────┐       │       │
             │       │       │          START          │       │       │
  ───────────┼───────┼───────┼────────────┼────────────┼───────┼───────┼──────────
             │       │       └────────────┼────────────┘       │       │
             │       └────────────────────┼────────────────────┘       │
             └────────────────────────────┼────────────────────────────┘
                                          │
            4. Plan the next phase        │      3. Develop & verify the
                                          │         next-level product
  Each loop = one phase; radius = cumulative cost; angle = progress through the quadrants
\end{Verbatim}
\end{listingblock}

\begin{itemize}
\tightlist
\item
  4 quadrants: Planning (objectives) → \textbf{Risk analysis} → Engineering → Evaluation. Meta-model; suits large, high-risk projects.
\end{itemize}

\needspace{130pt}

\hypertarget{p2-06--15-others}{%
\subsection{Others}\label{p2-06--15-others}}

\begingroup\small

\begin{longtable}[]{@{}
  >{\raggedright\arraybackslash}p{(\columnwidth - 2\tabcolsep) * \real{0.2450}}
  >{\raggedright\arraybackslash}p{(\columnwidth - 2\tabcolsep) * \real{0.7550}}@{}}
\toprule\noalign{}
\begin{minipage}[b]{\linewidth}\raggedright
\textbf{Model}
\end{minipage} & \begin{minipage}[b]{\linewidth}\raggedright
\textbf{Key idea}
\end{minipage} \\
\midrule\noalign{}
\endhead
\bottomrule\noalign{}
\endlastfoot
Concurrent & Activities exist simultaneously in different states (state chart) \\
Component-based (CBSE) & Assemble pre-built COTS components \\
Formal methods & Mathematical specification (Z, VDM); \textbf{Cleanroom} SE (statistical testing, defect prevention) \\
AOSD & Aspect-oriented: crosscutting concerns (logging, security) as aspects \\
\textbf{Unified Process (RUP)} & Use-case driven, architecture-centric, iterative \& incremental. Phases: \textbf{Inception → Elaboration → Construction → Transition} (+ Production) \\
\end{longtable}

\endgroup

\needspace{125pt}

\hypertarget{p2-06--16-agile}{%
\subsection{Agile}\label{p2-06--16-agile}}

\textbf{Agile Manifesto (2001)} values:

\begin{itemize}
\tightlist
\item
  \textbf{Individuals \& interactions} over processes \& tools
\item
  \textbf{Working software} over comprehensive documentation
\item
  \textbf{Customer collaboration} over contract negotiation
\item
  \textbf{Responding to change} over following a plan
  (12 principles; e.g. deliver frequently, welcome changing requirements, sustainable pace, face-to-face communication.)
\end{itemize}

\begingroup\small

\begin{longtable}[]{@{}
  >{\raggedright\arraybackslash}p{(\columnwidth - 2\tabcolsep) * \real{0.1583}}
  >{\raggedright\arraybackslash}p{(\columnwidth - 2\tabcolsep) * \real{0.8417}}@{}}
\toprule\noalign{}
\begin{minipage}[b]{\linewidth}\raggedright
\textbf{Method}
\end{minipage} & \begin{minipage}[b]{\linewidth}\raggedright
\textbf{Highlights}
\end{minipage} \\
\midrule\noalign{}
\endhead
\bottomrule\noalign{}
\endlastfoot
\textbf{XP} (Kent Beck) & Pair programming, \textbf{test-first (TDD)}, refactoring, continuous integration, small releases, \textbf{user stories}, spike solutions, CRC cards, collective ownership, 40-hour week \\
\textbf{Scrum} & \textbf{Sprint} (2--4 weeks, time-boxed), \textbf{Product backlog}, Sprint backlog, \textbf{Daily Scrum} (15 min stand-up), Sprint review, retrospective; roles: \textbf{Product Owner, Scrum Master, Development Team}; burndown chart \\
DSDM & Time-boxing, \textbf{MoSCoW} prioritisation (Must, Should, Could, Won\textquotesingle t); 80/20 principle \\
Crystal & Family (Clear, Orange\ldots) by team size \& criticality (Alistair Cockburn) \\
FDD & Feature Driven Development; features in ≤ 2 weeks \\
ASD & Adaptive SD: Speculation, Collaboration, Learning (Highsmith) \\
Kanban & Visual board, limit Work-In-Progress (WIP) \\
Lean & Eliminate waste \\
\end{longtable}

\endgroup

\diagram{figures/p2-06-01}

\needspace{130pt}

\hypertarget{p2-06--model-selection-cheat-sheet}{%
\subsection{Model selection cheat-sheet}\label{p2-06--model-selection-cheat-sheet}}

\begingroup\small

\begin{longtable}[]{@{}
  >{\raggedright\arraybackslash}p{(\columnwidth - 2\tabcolsep) * \real{0.3233}}
  >{\raggedright\arraybackslash}p{(\columnwidth - 2\tabcolsep) * \real{0.2633}}@{}}
\toprule\noalign{}
\begin{minipage}[b]{\linewidth}\raggedright
\textbf{Situation}
\end{minipage} & \begin{minipage}[b]{\linewidth}\raggedright
\textbf{Model}
\end{minipage} \\
\midrule\noalign{}
\endhead
\bottomrule\noalign{}
\endlastfoot
Clear, fixed requirements & Waterfall \\
Unclear requirements & Prototyping \\
High risk, large & Spiral \\
Quick delivery, modular & RAD / Incremental \\
Changing requirements, small teams & Agile \\
Safety-critical & Formal methods / Cleanroom \\
\end{longtable}

\endgroup

\needspace{95pt}

\hypertarget{p2-06--2-requirements-engineering}{%
\section{Requirements Engineering}\label{p2-06--2-requirements-engineering}}

\diagram{figures/p2-06-02}

\begin{itemize}
\tightlist
\item
  \textbf{Functional}: what the system does (services). \textbf{Non-functional}: quality constraints --- performance, security, reliability, usability, portability, maintainability (FURPS+).
\item
  Elicitation techniques: interviews, questionnaires, \textbf{JAD} (Joint Application Development), \textbf{QFD} (Quality Function Deployment: normal, expected, exciting requirements), brainstorming, use cases, observation, prototyping.
\item
  \textbf{SRS} (IEEE 830): characteristics --- correct, unambiguous, complete, consistent, ranked, \textbf{verifiable}, modifiable, traceable. SRS says \textbf{what}, not \textbf{how}.
\item
  Requirement traceability matrix. Validation: reviews, prototyping, test-case generation.
\item
  Analysis models: \textbf{DFD} (process), ER (data), \textbf{state diagrams} (behaviour), use case diagrams, data dictionary.
\end{itemize}

\begin{listingblock}{61pt}
\begin{Verbatim}[fontsize=\fontsize{8.4}{10.2}\selectfont]
 DFD notation (Yourdon/DeMarco)
 ┌──────┐  External entity   ( ○ )  Process (bubble)   ═══  Data store   ──►  Data flow
 Level 0 DFD = Context diagram (single process, shows external entities only)
 Balancing: inputs/outputs of child DFD = those of parent process
\end{Verbatim}
\end{listingblock}

\needspace{95pt}

\hypertarget{p2-06--uml-diagrams}{%
\subsection{UML Diagrams}\label{p2-06--uml-diagrams}}

\diagram{figures/p2-06-03}

Relationships: association, aggregation (◇ hollow, "has-a", weak), composition (◆ filled, strong ownership), generalisation (△ "is-a"), dependency (dashed arrow), realisation. Use case relations: «include» (mandatory), «extend» (optional).

\needspace{249pt}

\hypertarget{p2-06--3-software-design}{%
\section{Software Design}\label{p2-06--3-software-design}}

\needspace{213pt}

\hypertarget{p2-06--31-design-concepts}{%
\subsection{Design Concepts}\label{p2-06--31-design-concepts}}

Abstraction (procedural, data), architecture, patterns, separation of concerns, \textbf{modularity}, \textbf{information hiding} (Parnas), \textbf{functional independence} (high cohesion + low coupling), stepwise refinement (Wirth), refactoring, aspects.

\needspace{147pt}

\hypertarget{p2-06--32-cohesion-within-a-module--best-to-worst}{%
\subsection{Cohesion (within a module) --- best to worst}\label{p2-06--32-cohesion-within-a-module--best-to-worst}}

\begin{listingblock}{97pt}
\begin{Verbatim}[fontsize=\fontsize{9.0}{10.9}\selectfont]
 BEST  ▲  Functional      – single well-defined task
       │  Sequential      – output of one part is input to next
       │  Communicational – parts operate on same data
       │  Procedural      – parts follow a sequence of execution
       │  Temporal        – executed at same time (e.g. initialisation)
       │  Logical         – logically similar tasks, selected by flag
 WORST ▼  Coincidental    – unrelated tasks
\end{Verbatim}
\end{listingblock}

Mnemonic (worst→best): \textbf{"Coin Logic Temp Proc Comm Seq Func"}.

\needspace{136pt}

\hypertarget{p2-06--33-coupling-between-modules--best-to-worst}{%
\subsection{Coupling (between modules) --- best to worst}\label{p2-06--33-coupling-between-modules--best-to-worst}}

\begin{listingblock}{86pt}
\begin{Verbatim}[fontsize=\fontsize{9.0}{10.9}\selectfont]
 BEST  ▲  No coupling / Data coupling – pass only needed data
       │  Stamp coupling   – pass whole data structure, use part
       │  Control coupling – pass control flag
       │  External coupling – share external format/device/protocol
       │  Common coupling  – share global data
 WORST ▼  Content coupling – one module modifies/uses internals of another
\end{Verbatim}
\end{listingblock}

\textbf{Goal: high cohesion, low coupling.}

\needspace{196pt}

\hypertarget{p2-06--34-architectural-styles}{%
\subsection{Architectural Styles}\label{p2-06--34-architectural-styles}}

Data-centred (repository/blackboard), data-flow (\textbf{pipe \& filter}, batch sequential), call-and-return (main/subprogram, remote procedure), object-oriented, layered, client-server, \textbf{MVC}, event-driven, microservices, SOA.

\needspace{130pt}

\hypertarget{p2-06--35-design-patterns-gof--23}{%
\subsection{Design Patterns (GoF -- 23)}\label{p2-06--35-design-patterns-gof--23}}

\begingroup\small

\begin{longtable}[]{@{}
  >{\raggedright\arraybackslash}p{(\columnwidth - 2\tabcolsep) * \real{0.1617}}
  >{\raggedright\arraybackslash}p{(\columnwidth - 2\tabcolsep) * \real{0.8383}}@{}}
\toprule\noalign{}
\begin{minipage}[b]{\linewidth}\raggedright
\textbf{Category}
\end{minipage} & \begin{minipage}[b]{\linewidth}\raggedright
\textbf{Patterns}
\end{minipage} \\
\midrule\noalign{}
\endhead
\bottomrule\noalign{}
\endlastfoot
Creational (5) & Singleton, Factory Method, Abstract Factory, Builder, Prototype \\
Structural (7) & Adapter, Bridge, Composite, Decorator, Facade, Flyweight, Proxy \\
Behavioural (11) & Observer, Strategy, Command, Iterator, State, Template Method, Visitor, Mediator, Memento, Chain of Responsibility, Interpreter \\
\end{longtable}

\endgroup

\needspace{197pt}

\hypertarget{p2-06--36-ui-design}{%
\subsection{UI Design}\label{p2-06--36-ui-design}}

Golden rules (Mandel): place user in control, reduce memory load, make interface consistent. Usability heuristics (Nielsen).

\needspace{131pt}

\hypertarget{p2-06--4-software-quality}{%
\section{Software Quality}\label{p2-06--4-software-quality}}

\needspace{95pt}

\hypertarget{p2-06--mccalls-quality-factors-1977--11-factors-in-3-perspectives}{%
\subsection{McCall\textquotesingle s Quality Factors (1977) --- 11 factors in 3 perspectives}\label{p2-06--mccalls-quality-factors-1977--11-factors-in-3-perspectives}}

\diagram{figures/p2-06-04}

\textbf{ISO 9126} (6 characteristics): \textbf{Functionality, Reliability, Usability, Efficiency, Maintainability, Portability} (FRUEMP). Replaced by \textbf{ISO/IEC 25010} (8: + Security, Compatibility; "functional suitability", "performance efficiency").

\begin{itemize}
\tightlist
\item
  \textbf{Quality control (QC)}: inspections, reviews, tests --- product-oriented, detect defects.
\item
  \textbf{Quality assurance (QA)}: auditing \& reporting --- process-oriented, prevent defects.
\item
  \textbf{SQA} activities: FTR (Formal Technical Reviews --- walkthrough, inspection: 3--5 people, ≤ 2 hours prep, ≤ 2 hours meeting), standards, audits.
\item
  \textbf{Six Sigma}: 3.4 defects per million opportunities; DMAIC (Define, Measure, Analyse, Improve, Control).
\item
  \textbf{CMMI levels}: 1 Initial → 2 Managed (Repeatable) → 3 Defined → 4 Quantitatively Managed → \textbf{5 Optimising}.
\item
  ISO 9001 --- QMS standard applicable to software.
\end{itemize}

\needspace{114pt}

\hypertarget{p2-06--reliability}{%
\subsection{Reliability}\label{p2-06--reliability}}

\begin{listingblock}{64pt}
\begin{Verbatim}[fontsize=\fontsize{9.0}{10.9}\selectfont]
MTBF = MTTF + MTTR
Availability = MTTF / (MTTF + MTTR) × 100%
ROCOF = rate of occurrence of failure; POFOD = probability of failure on demand
Reliability R(t) = e^(−λt)   (λ = failure rate)
\end{Verbatim}
\end{listingblock}

\needspace{95pt}

\hypertarget{p2-06--risk-management}{%
\subsection{Risk Management}\label{p2-06--risk-management}}

\begin{itemize}
\tightlist
\item
  \textbf{Reactive} (fire-fighting) vs \textbf{Proactive}.
\item
  Risk categories: project, technical, business; known, predictable, unpredictable.
\item
  \textbf{Risk exposure RE = P × C} (probability × cost/impact).
\item
  Risk table → \textbf{RMMM}: Risk \textbf{Mitigation} (avoid), \textbf{Monitoring}, \textbf{Management} (contingency).
\end{itemize}

\needspace{157pt}

\hypertarget{p2-06--5-estimation}{%
\section{Estimation}\label{p2-06--5-estimation}}

\needspace{121pt}

\hypertarget{p2-06--51-loc--function-points}{%
\subsection{LOC \& Function Points}\label{p2-06--51-loc--function-points}}

\textbf{Function Point (Albrecht, 1979)}:

\begin{listingblock}{41pt}
\begin{Verbatim}[fontsize=\fontsize{8.6}{10.4}\selectfont]
FP = UFP × VAF        VAF = 0.65 + 0.01 × ΣFᵢ     (14 GSCs, each 0–5 → ΣFᵢ ∈ [0, 70])
VAF range: 0.65 to 1.35
\end{Verbatim}
\end{listingblock}

\begingroup\small

\begin{longtable}[]{@{}
  >{\raggedright\arraybackslash}p{(\columnwidth - 6\tabcolsep) * \real{0.3196}}
  >{\raggedright\arraybackslash}p{(\columnwidth - 6\tabcolsep) * \real{0.0932}}
  >{\raggedright\arraybackslash}p{(\columnwidth - 6\tabcolsep) * \real{0.0979}}
  >{\raggedright\arraybackslash}p{(\columnwidth - 6\tabcolsep) * \real{0.1059}}@{}}
\toprule\noalign{}
\begin{minipage}[b]{\linewidth}\raggedright
\textbf{Parameter}
\end{minipage} & \begin{minipage}[b]{\linewidth}\raggedright
\textbf{Simple}
\end{minipage} & \begin{minipage}[b]{\linewidth}\raggedright
\textbf{Average}
\end{minipage} & \begin{minipage}[b]{\linewidth}\raggedright
\textbf{Complex}
\end{minipage} \\
\midrule\noalign{}
\endhead
\bottomrule\noalign{}
\endlastfoot
External Inputs (EI) & 3 & 4 & 6 \\
External Outputs (EO) & 4 & 5 & 7 \\
External Inquiries (EQ) & 3 & 4 & 6 \\
Internal Logical Files (ILF) & 7 & 10 & 15 \\
External Interface Files (EIF) & 5 & 7 & 10 \\
\end{longtable}

\endgroup

\textbf{Worked}: EI = 10 (avg), EO = 8 (avg), EQ = 5 (avg), ILF = 4 (avg), EIF = 2 (avg); ΣFᵢ = 42.

\begin{listingblock}{53pt}
\begin{Verbatim}[fontsize=\fontsize{9.0}{10.9}\selectfont]
UFP = 10×4 + 8×5 + 5×4 + 4×10 + 2×7 = 40 + 40 + 20 + 40 + 14 = 154
VAF = 0.65 + 0.42 = 1.07
FP  = 154 × 1.07 = 164.78
\end{Verbatim}
\end{listingblock}

\needspace{103pt}

\hypertarget{p2-06--52-cocomo-boehm-1981}{%
\subsection{COCOMO (Boehm, 1981)}\label{p2-06--52-cocomo-boehm-1981}}

\begin{listingblock}{53pt}
\begin{Verbatim}[fontsize=\fontsize{9.0}{10.9}\selectfont]
Basic COCOMO:  Effort E = a·(KLOC)^b  person-months
               Time   D = c·(E)^d     months
               Staff  = E / D
\end{Verbatim}
\end{listingblock}

\begingroup\small

\begin{longtable}[]{@{}
  >{\raggedright\arraybackslash}p{(\columnwidth - 10\tabcolsep) * \real{0.1635}}
  >{\raggedright\arraybackslash}p{(\columnwidth - 10\tabcolsep) * \real{0.0551}}
  >{\raggedright\arraybackslash}p{(\columnwidth - 10\tabcolsep) * \real{0.0679}}
  >{\raggedright\arraybackslash}p{(\columnwidth - 10\tabcolsep) * \real{0.0551}}
  >{\raggedright\arraybackslash}p{(\columnwidth - 10\tabcolsep) * \real{0.0679}}
  >{\raggedright\arraybackslash}p{(\columnwidth - 10\tabcolsep) * \real{0.3692}}@{}}
\toprule\noalign{}
\begin{minipage}[b]{\linewidth}\raggedright
\textbf{Mode}
\end{minipage} & \begin{minipage}[b]{\linewidth}\raggedright
\textbf{a}
\end{minipage} & \begin{minipage}[b]{\linewidth}\raggedright
\textbf{b}
\end{minipage} & \begin{minipage}[b]{\linewidth}\raggedright
\textbf{c}
\end{minipage} & \begin{minipage}[b]{\linewidth}\raggedright
\textbf{d}
\end{minipage} & \begin{minipage}[b]{\linewidth}\raggedright
\textbf{Description}
\end{minipage} \\
\midrule\noalign{}
\endhead
\bottomrule\noalign{}
\endlastfoot
\textbf{Organic} & 2.4 & 1.05 & 2.5 & 0.38 & Small team, familiar, \textless{} 50 KLOC \\
\textbf{Semi-detached} & 3.0 & 1.12 & 2.5 & 0.35 & Medium, mixed experience \\
\textbf{Embedded} & 3.6 & 1.20 & 2.5 & 0.32 & Tight constraints, hardware \\
\end{longtable}

\endgroup

\textbf{Worked}: Organic, 32 KLOC.

\begin{listingblock}{53pt}
\begin{Verbatim}[fontsize=\fontsize{9.0}{10.9}\selectfont]
E = 2.4 × 32^1.05 ≈ 2.4 × 38.06 ≈ 91 PM
D = 2.5 × 91^0.38 ≈ 2.5 × 5.55 ≈ 13.9 months
Staff ≈ 91 / 13.9 ≈ 6.5 persons
\end{Verbatim}
\end{listingblock}

\begin{itemize}
\tightlist
\item
  \textbf{Intermediate COCOMO}: E = a·KLOC\^{}b × \textbf{EAF} (product of \textbf{15 cost drivers} in 4 groups: product, hardware/computer, personnel, project). Coefficients a: 3.2, 3.0, 2.8.
\item
  \textbf{Detailed COCOMO}: phase-sensitive effort multipliers.
\item
  \textbf{COCOMO II}: Application composition (object points), Early design, Post-architecture models; scale factors.
\item
  \textbf{Putnam/SLIM} model: Rayleigh curve; L = Ck · K\^{}(1/3) · td\^{}(4/3).
\item
  \textbf{Brooks\textquotesingle{} law}: "Adding manpower to a late project makes it later." Communication paths among n people = n(n−1)/2.
\end{itemize}

\needspace{95pt}

\hypertarget{p2-06--53-scheduling}{%
\subsection{Scheduling}\label{p2-06--53-scheduling}}

\begin{itemize}
\tightlist
\item
  \textbf{WBS} (work breakdown structure), \textbf{Gantt / timeline charts}, \textbf{PERT/CPM} (critical path; see \protect\hyperlink{p2-01--75-pert-cpm}{Unit 1}).
\item
  \textbf{Earned Value Analysis}: BCWS (planned value), BCWP (earned value), ACWP (actual cost).

  \begin{itemize}
  \tightlist
  \item
    SPI = BCWP/BCWS; CPI = BCWP/ACWP; SV = BCWP − BCWS; CV = BCWP − ACWP. SPI \textless{} 1 → behind schedule.
  \end{itemize}
\item
  40-20-40 rule: 40\% analysis \& design, 20\% coding, 40\% testing.
\end{itemize}

\needspace{116pt}

\hypertarget{p2-06--6-software-testing}{%
\section{Software Testing}\label{p2-06--6-software-testing}}

\needspace{80pt}

\hypertarget{p2-06--61-terminology}{%
\subsection{Terminology}\label{p2-06--61-terminology}}

\begin{listingblock}{30pt}
\begin{Verbatim}[fontsize=\fontsize{7.5}{9.1}\selectfont]
Mistake/Error (human) ──► Fault/Defect/Bug (in code) ──► Failure (observed deviation at run time)
\end{Verbatim}
\end{listingblock}

\begin{itemize}
\tightlist
\item
  \textbf{Verification}: "Are we building the product right?" (static --- reviews, inspections).
\item
  \textbf{Validation}: "Are we building the right product?" (dynamic --- testing against user needs).
\item
  Testing shows the \textbf{presence} of bugs, not their absence (Dijkstra).
\end{itemize}

\needspace{95pt}

\hypertarget{p2-06--62-levels-of-testing}{%
\subsection{Levels of Testing}\label{p2-06--62-levels-of-testing}}

\diagram{figures/p2-06-05}

\begin{itemize}
\tightlist
\item
  \textbf{Stub} = dummy called module (top-down). \textbf{Driver} = dummy calling module (bottom-up).
\item
  \textbf{Regression testing}: re-run tests after changes. \textbf{Smoke testing}: build verification. \textbf{Sanity}: narrow check after fix.
\item
  \textbf{Alpha}: by customers at developer\textquotesingle s site (controlled). \textbf{Beta}: by end users at their site (uncontrolled).
\item
  Stress (beyond limits), load (expected load), volume (large data), performance, recovery, security testing.
\end{itemize}

\needspace{130pt}

\hypertarget{p2-06--63-black-box-vs-white-box}{%
\subsection{Black-box vs White-box}\label{p2-06--63-black-box-vs-white-box}}

\begingroup\small

\begin{longtable}[]{@{}
  >{\raggedright\arraybackslash}p{(\columnwidth - 2\tabcolsep) * \real{0.3508}}
  >{\raggedright\arraybackslash}p{(\columnwidth - 2\tabcolsep) * \real{0.3800}}@{}}
\toprule\noalign{}
\begin{minipage}[b]{\linewidth}\raggedright
\textbf{Black-box (functional, behavioural)}
\end{minipage} & \begin{minipage}[b]{\linewidth}\raggedright
\textbf{White-box (structural, glass-box)}
\end{minipage} \\
\midrule\noalign{}
\endhead
\bottomrule\noalign{}
\endlastfoot
Equivalence partitioning & Statement coverage \\
\textbf{Boundary value analysis} & Branch/\allowbreak{}decision coverage \\
Cause-effect graphing / decision tables & Condition, MC/DC coverage \\
State transition testing & \textbf{Basis path testing} (cyclomatic complexity) \\
Orthogonal array testing & Data-flow testing, loop testing \\
Error guessing & Mutation testing \\
\end{longtable}

\endgroup

\begin{itemize}
\tightlist
\item
  \textbf{Boundary value analysis} for range {[}1, 100{]}: test 0, 1, 2, 99, 100, 101 (robust) or 1, 2, 50, 99, 100 (normal: 4n + 1 for n variables).
\item
  \textbf{Equivalence classes} for {[}1, 100{]}: one valid (1--100), two invalid (\textless{} 1, \textgreater{} 100).
\item
  Coverage strength: Statement \textless{} Branch \textless{} Condition/Decision \textless{} MC/DC \textless{} Multiple condition \textless{} Path.
\end{itemize}

\needspace{111pt}

\hypertarget{p2-06--64-cyclomatic-complexity-mccabe-1976}{%
\subsection{Cyclomatic Complexity (McCabe, 1976)}\label{p2-06--64-cyclomatic-complexity-mccabe-1976}}

\begin{listingblock}{61pt}
\begin{Verbatim}[fontsize=\fontsize{8.4}{10.2}\selectfont]
V(G) = E − N + 2P        (P = connected components, usually 1)
V(G) = number of predicate (decision) nodes + 1
V(G) = number of bounded regions + 1  (regions of planar flow graph incl. outer)
     = number of linearly independent paths = minimum test cases for basis path testing
\end{Verbatim}
\end{listingblock}

\begin{listingblock}{140pt}
\begin{Verbatim}[fontsize=\fontsize{9.0}{10.9}\selectfont]
 Flow graph for: if (a) { x } else { y }; while (b) { z }
      (1) if a
      ╱    ╲
    (2)x   (3)y
      ╲    ╱
      (4) while b ◄──┐
        │   ╲        │
        │   (5) z ───┘
        ▼
       (6) end
 N = 6, E = 7  → V(G) = 7 − 6 + 2 = 3   (2 decisions + 1 = 3)
\end{Verbatim}
\end{listingblock}

V(G) ≤ 10 recommended.

\needspace{95pt}

\hypertarget{p2-06--7-software-configuration-management}{%
\section{Software Configuration Management}\label{p2-06--7-software-configuration-management}}

\begin{itemize}
\tightlist
\item
  \textbf{SCM} activities: identification, version control, change control, configuration auditing, status reporting.
\item
  \textbf{Baseline}: formally reviewed \& agreed work product; changes only via formal change control (Change Control Board / Authority).
\item
  SCIs (Software Configuration Items). Tools: Git, SVN, CVS.
\end{itemize}

\needspace{130pt}

\hypertarget{p2-06--8-maintenance--re-engineering}{%
\section{Maintenance \& Re-engineering}\label{p2-06--8-maintenance--re-engineering}}

\begingroup\small

\begin{longtable}[]{@{}
  >{\raggedright\arraybackslash}p{(\columnwidth - 4\tabcolsep) * \real{0.1801}}
  >{\raggedright\arraybackslash}p{(\columnwidth - 4\tabcolsep) * \real{0.1717}}
  >{\raggedright\arraybackslash}p{(\columnwidth - 4\tabcolsep) * \real{0.4057}}@{}}
\toprule\noalign{}
\begin{minipage}[b]{\linewidth}\raggedright
\textbf{Maintenance type}
\end{minipage} & \begin{minipage}[b]{\linewidth}\raggedright
\textbf{Share (approx.)}
\end{minipage} & \begin{minipage}[b]{\linewidth}\raggedright
\textbf{Purpose}
\end{minipage} \\
\midrule\noalign{}
\endhead
\bottomrule\noalign{}
\endlastfoot
Corrective & \textasciitilde20\% & Fix bugs \\
\textbf{Adaptive} & \textasciitilde25\% & Adapt to new environment (OS, hardware) \\
\textbf{Perfective} & \textasciitilde50\% (largest) & New features, performance \\
Preventive & \textasciitilde5\% & Improve future maintainability (refactoring) \\
\end{longtable}

\endgroup

\begin{itemize}
\tightlist
\item
  \textbf{Lehman\textquotesingle s laws} of software evolution (continuing change, increasing complexity\ldots).
\item
  \textbf{Re-engineering}: inventory analysis → document restructuring → \textbf{reverse engineering} → code restructuring → data restructuring → \textbf{forward engineering}.
\item
  \textbf{Reverse engineering}: extracting design/specification from code (abstraction level ↑). Forward engineering: design → code.
\item
  \textbf{Restructuring}: change form without changing functionality.
\item
  Software maintenance cost is typically \textbf{60--80\%} of total lifecycle cost.
\end{itemize}

\needspace{127pt}

\hypertarget{p2-06--other-metrics}{%
\subsection{Other Metrics}\label{p2-06--other-metrics}}

\begin{listingblock}{77pt}
\begin{Verbatim}[fontsize=\fontsize{7.8}{9.5}\selectfont]
Halstead:  n1 = distinct operators, n2 = distinct operands, N1, N2 = totals
           Vocabulary n = n1 + n2 ; Length N = N1 + N2
           Volume V = N log₂ n ; Difficulty D = (n1/2)(N2/n2) ; Effort E = D × V
Defect density = defects / KLOC
DRE (Defect Removal Efficiency) = E / (E + D)   E = errors before delivery, D = defects after
CK (OO) metrics: WMC, DIT, NOC, CBO, RFC, LCOM
\end{Verbatim}
\end{listingblock}

\needspace{131pt}

\hypertarget{p2-06--9-deeper-dive--modelling-examples-test-design-estimation--reliability-models}{%
\section{Deeper Dive --- Modelling Examples, Test Design, Estimation \& Reliability Models}\label{p2-06--9-deeper-dive--modelling-examples-test-design-estimation--reliability-models}}

\needspace{95pt}

\hypertarget{p2-06--91-level-1-dfd--library-system}{%
\subsection{Level-1 DFD --- Library System}\label{p2-06--91-level-1-dfd--library-system}}

\diagram{figures/p2-06-06}

Rules: every process has at least one input and one output; data stores connect only through processes; external entities never connect directly to data stores.

\needspace{95pt}

\hypertarget{p2-06--92-uml-class-diagram}{%
\subsection{UML Class Diagram}\label{p2-06--92-uml-class-diagram}}

\diagram{figures/p2-06-07}

Visibility: \texttt{+} public, \texttt{-} private, \texttt{\#} protected, \texttt{\textasciitilde{}} package. Multiplicity: 1, 0..1, 0..\emph{, 1..}.

\needspace{95pt}

\hypertarget{p2-06--93-uml-sequence-diagram--atm-withdrawal}{%
\subsection{UML Sequence Diagram --- ATM Withdrawal}\label{p2-06--93-uml-sequence-diagram--atm-withdrawal}}

\diagram{figures/p2-06-08}

\needspace{160pt}

\hypertarget{p2-06--94-black-box-test-design--worked}{%
\subsection{Black-box Test Design --- Worked}\label{p2-06--94-black-box-test-design--worked}}

Function: \texttt{eligible(age,\ grade)} where age ∈ {[}18, 60{]} and grade ∈ \{A, B, C\}.

\begingroup\small

\begin{longtable}[]{@{}
  >{\raggedright\arraybackslash}p{(\columnwidth - 2\tabcolsep) * \real{0.2756}}
  >{\raggedright\arraybackslash}p{(\columnwidth - 2\tabcolsep) * \real{0.6172}}@{}}
\toprule\noalign{}
\begin{minipage}[b]{\linewidth}\raggedright
\textbf{Technique}
\end{minipage} & \begin{minipage}[b]{\linewidth}\raggedright
\textbf{Test values}
\end{minipage} \\
\midrule\noalign{}
\endhead
\bottomrule\noalign{}
\endlastfoot
Equivalence classes (age) & valid 18--60 (e.g. 35); invalid \textless{} 18 (e.g. 10); invalid \textgreater{} 60 (e.g. 70) \\
Equivalence classes (grade) & valid \{A, B, C\}; invalid (e.g. D) \\
BVA (age, normal) & 18, 19, 39, 59, 60 \\
Robust BVA (age) & 17, 18, 19, 39, 59, 60, 61 \\
Worst-case BVA (2 variables) & 5² = 25 combinations; robust worst-case 7² = 49 \\
\end{longtable}

\endgroup

\textbf{Decision table} (login):

\begingroup\small

\begin{longtable}[]{@{}
  >{\raggedright\arraybackslash}p{(\columnwidth - 8\tabcolsep) * \real{0.2451}}
  >{\raggedright\arraybackslash}p{(\columnwidth - 8\tabcolsep) * \real{0.0510}}
  >{\raggedright\arraybackslash}p{(\columnwidth - 8\tabcolsep) * \real{0.0510}}
  >{\raggedright\arraybackslash}p{(\columnwidth - 8\tabcolsep) * \real{0.0510}}
  >{\raggedright\arraybackslash}p{(\columnwidth - 8\tabcolsep) * \real{0.0510}}@{}}
\toprule\noalign{}
\begin{minipage}[b]{\linewidth}\raggedright
\textbf{Conditions / Actions}
\end{minipage} & \begin{minipage}[b]{\linewidth}\raggedright
\textbf{R1}
\end{minipage} & \begin{minipage}[b]{\linewidth}\raggedright
\textbf{R2}
\end{minipage} & \begin{minipage}[b]{\linewidth}\raggedright
\textbf{R3}
\end{minipage} & \begin{minipage}[b]{\linewidth}\raggedright
\textbf{R4}
\end{minipage} \\
\midrule\noalign{}
\endhead
\bottomrule\noalign{}
\endlastfoot
Valid user ID? & T & T & F & F \\
Valid password? & T & F & T & F \\
Grant access & ✔ & & & \\
Show "wrong password" & & ✔ & & \\
Show "unknown user" & & & ✔ & ✔ \\
\end{longtable}

\endgroup

\needspace{140pt}

\hypertarget{p2-06--95-intermediate-cocomo--worked}{%
\subsection{Intermediate COCOMO --- Worked}\label{p2-06--95-intermediate-cocomo--worked}}

Semi-detached project, 50 KLOC, EAF = 1.2 (product of the 15 cost-driver multipliers).

\begin{listingblock}{60pt}
\begin{Verbatim}[fontsize=\fontsize{8.2}{10.0}\selectfont]
Intermediate coefficients a: organic 3.2, semi-detached 3.0, embedded 2.8 (b as in basic)
E = 3.0 × 50^1.12 × 1.2 ≈ 3.0 × 80 × 1.2 ≈ 288 person-months
D = 2.5 × 288^0.35 ≈ 2.5 × 7.26 ≈ 18 months
Average staff ≈ 288 / 18 = 16 persons
\end{Verbatim}
\end{listingblock}

Cost drivers include RELY (reliability), CPLX (complexity), ACAP (analyst capability), PCAP, TOOL, SCED (schedule constraint), etc.

\needspace{130pt}

\hypertarget{p2-06--96-software-reliability-models}{%
\subsection{Software Reliability Models}\label{p2-06--96-software-reliability-models}}

\begingroup\small

\begin{longtable}[]{@{}
  >{\raggedright\arraybackslash}p{(\columnwidth - 2\tabcolsep) * \real{0.2483}}
  >{\raggedright\arraybackslash}p{(\columnwidth - 2\tabcolsep) * \real{0.7094}}@{}}
\toprule\noalign{}
\begin{minipage}[b]{\linewidth}\raggedright
\textbf{Model}
\end{minipage} & \begin{minipage}[b]{\linewidth}\raggedright
\textbf{Idea}
\end{minipage} \\
\midrule\noalign{}
\endhead
\bottomrule\noalign{}
\endlastfoot
\textbf{Jelinski--Moranda} (1972) & N initial faults; each fix reduces failure rate by a constant φ: λᵢ = φ(N − i + 1) \\
\textbf{Musa basic execution time} & λ(μ) = λ₀ (1 − μ/ν₀): failure intensity falls linearly with failures experienced \\
Musa--Okumoto logarithmic & Failure intensity decreases exponentially with failures experienced \\
Goel--Okumoto (NHPP) & Expected failures m(t) = a(1 − e\^{}(−bt)) \\
Littlewood--Verrall & Bayesian; failure rates random \\
\end{longtable}

\endgroup

\textbf{Worked (Musa basic)}: λ₀ = 10 failures/CPU-hr, ν₀ = 100 total failures. After 50 failures: λ = 10(1 − 50/100) = \textbf{5 failures/CPU-hr}.

\needspace{130pt}

\hypertarget{p2-06--97-risk-table--example}{%
\subsection{Risk Table --- Example}\label{p2-06--97-risk-table--example}}

\begingroup\small

\begin{longtable}[]{@{}
  >{\raggedright\arraybackslash}p{(\columnwidth - 10\tabcolsep) * \real{0.1861}}
  >{\raggedright\arraybackslash}p{(\columnwidth - 10\tabcolsep) * \real{0.1147}}
  >{\raggedright\arraybackslash}p{(\columnwidth - 10\tabcolsep) * \real{0.1427}}
  >{\raggedright\arraybackslash}p{(\columnwidth - 10\tabcolsep) * \real{0.1808}}
  >{\raggedright\arraybackslash}p{(\columnwidth - 10\tabcolsep) * \real{0.1463}}
  >{\raggedright\arraybackslash}p{(\columnwidth - 10\tabcolsep) * \real{0.2294}}@{}}
\toprule\noalign{}
\begin{minipage}[b]{\linewidth}\raggedright
\textbf{Risk}
\end{minipage} & \begin{minipage}[b]{\linewidth}\raggedright
\textbf{Category}
\end{minipage} & \begin{minipage}[b]{\linewidth}\raggedright
\textbf{Probability}
\end{minipage} & \begin{minipage}[b]{\linewidth}\raggedright
\textbf{Impact (1--4)}
\end{minipage} & \begin{minipage}[b]{\linewidth}\raggedright
\textbf{Exposure (₹)}
\end{minipage} & \begin{minipage}[b]{\linewidth}\raggedright
\textbf{RMMM response}
\end{minipage} \\
\midrule\noalign{}
\endhead
\bottomrule\noalign{}
\endlastfoot
Key developer leaves & Project & 0.3 & 2 & 0.3 × 4 L = 1.2 L & Cross-training, documentation \\
Requirements change late & Business & 0.6 & 2 & 0.6 × 2 L = 1.2 L & Agile iterations, change control \\
New tool unreliable & Technical & 0.2 & 3 & 0.2 × 1 L = 0.2 L & Prototype tool early \\
\end{longtable}

\endgroup

Risks are sorted by probability × impact; a \textbf{cut-off line} decides which get full RMMM plans.

\needspace{130pt}

\hypertarget{p2-06--98-cmmi-process-areas-by-level-examples}{%
\subsection{CMMI Process Areas by Level (examples)}\label{p2-06--98-cmmi-process-areas-by-level-examples}}

\begingroup\small

\begin{longtable}[]{@{}
  >{\raggedright\arraybackslash}p{(\columnwidth - 4\tabcolsep) * \real{0.2374}}
  >{\raggedright\arraybackslash}p{(\columnwidth - 4\tabcolsep) * \real{0.3236}}
  >{\raggedright\arraybackslash}p{(\columnwidth - 4\tabcolsep) * \real{0.4391}}@{}}
\toprule\noalign{}
\begin{minipage}[b]{\linewidth}\raggedright
\textbf{Level}
\end{minipage} & \begin{minipage}[b]{\linewidth}\raggedright
\textbf{Focus}
\end{minipage} & \begin{minipage}[b]{\linewidth}\raggedright
\textbf{Example process areas}
\end{minipage} \\
\midrule\noalign{}
\endhead
\bottomrule\noalign{}
\endlastfoot
2 Managed & Basic project management & Requirements management, project planning, configuration management, measurement \& analysis, QA \\
3 Defined & Organisation-wide standard process & Requirements development, technical solution, verification, validation, risk management \\
4 Quantitatively managed & Statistical control & Organisational process performance, quantitative project management \\
5 Optimising & Continuous improvement & Causal analysis \& resolution, organisational performance management \\
\end{longtable}

\endgroup

\needspace{125pt}

\hypertarget{p2-06--previous-year-questions-pyq-pattern}{%
\section{Previous Year Questions (PYQ pattern)}\label{p2-06--previous-year-questions-pyq-pattern}}

\textbf{Process models}

\begin{enumerate}
\def\labelenumi{\arabic{enumi}.}
\tightlist
\item
  Which model is most suitable when requirements are not clear? \textbf{Prototyping}
\item
  The spiral model was proposed by: \textbf{Barry Boehm}; its key feature: \textbf{risk analysis}
\item
  In the spiral model, the radial dimension represents: \textbf{cumulative cost}
\item
  RUP phases in order: \textbf{Inception, Elaboration, Construction, Transition}
\item
  Which is NOT an agile method? (a) XP (b) Scrum (c) Waterfall (d) DSDM --- \textbf{Ans: (c)}
\item
  Pair programming is a practice of: \textbf{Extreme Programming}
\item
  In Scrum, the time-boxed iteration is called: \textbf{Sprint}; daily meeting lasts: \textbf{15 minutes}
\item
  MoSCoW prioritisation is associated with: \textbf{DSDM}
\item
  V-model emphasises: \textbf{verification and validation at each phase}
\end{enumerate}

\textbf{Requirements \& design}

\begin{enumerate}
\def\labelenumi{\arabic{enumi}.}
\setcounter{enumi}{9}
\tightlist
\item
  A level-0 DFD is also called: \textbf{context diagram}
\item
  "The system shall respond within 2 seconds" is a: \textbf{non-functional requirement}
\item
  SRS should specify: \textbf{what the system should do, not how}
\item
  Best type of cohesion: \textbf{functional}; worst: \textbf{coincidental}
\item
  Best coupling: \textbf{data coupling}; worst: \textbf{content coupling}
\item
  Modules sharing global data exhibit: \textbf{common coupling}
\item
  Passing a flag to control another module\textquotesingle s logic: \textbf{control coupling}
\item
  Information hiding was proposed by: \textbf{David Parnas}
\item
  Which UML diagram shows the time-ordered interaction of objects? \textbf{Sequence diagram}
\item
  «include» vs «extend»: include is \textbf{mandatory}, extend is \textbf{optional/conditional}
\item
  Singleton pattern belongs to: \textbf{creational patterns}
\item
  Filled diamond in UML indicates: \textbf{composition}
\end{enumerate}

\textbf{Quality}

\begin{enumerate}
\def\labelenumi{\arabic{enumi}.}
\setcounter{enumi}{21}
\tightlist
\item
  Which is NOT a McCall product-operation factor? (a) correctness (b) reliability (c) portability (d) usability --- \textbf{Ans: (c)}
\item
  ISO 9126 quality characteristics count: \textbf{6}
\item
  CMMI level 5 is: \textbf{Optimising}
\item
  Availability with MTTF = 90 h and MTTR = 10 h: \textbf{90\%}
\item
  Risk exposure with probability 0.3 and cost ₹50,000: \textbf{₹15,000}
\item
  RMMM stands for: \textbf{Risk Mitigation, Monitoring and Management}
\item
  QA is \_\_\_\_-oriented while QC is \_\_\_\_-oriented: \textbf{process; product}
\end{enumerate}

\textbf{Estimation}

\begin{enumerate}
\def\labelenumi{\arabic{enumi}.}
\setcounter{enumi}{28}
\tightlist
\item
  Value adjustment factor range: \textbf{0.65 to 1.35}
\item
  Number of general system characteristics in FP: \textbf{14}
\item
  Basic COCOMO organic effort constants: \textbf{a = 2.4, b = 1.05}
\item
  Intermediate COCOMO uses how many cost drivers? \textbf{15}
\item
  "Adding manpower to a late project makes it later" --- \textbf{Brooks\textquotesingle{} law}
\item
  Number of communication paths among 6 team members: \textbf{15}
\item
  SPI = 0.8 indicates the project is: \textbf{behind schedule}
\end{enumerate}

\textbf{Testing}

\begin{enumerate}
\def\labelenumi{\arabic{enumi}.}
\setcounter{enumi}{35}
\tightlist
\item
  Cyclomatic complexity of a graph with 10 edges, 8 nodes: 10 − 8 + 2 = \textbf{4}
\item
  A program has 3 \texttt{if} and 1 \texttt{while}: V(G) = 4 + 1 = \textbf{5}
\item
  Boundary value testing is a: \textbf{black-box technique}
\item
  Basis path testing is: \textbf{white-box}
\item
  Testing at developer\textquotesingle s site by customer: \textbf{alpha testing}
\item
  Re-running tests after modification: \textbf{regression testing}
\item
  In top-down integration, dummy modules are: \textbf{stubs}
\item
  "Are we building the product right?" --- \textbf{Verification}
\item
  Mutation testing evaluates: \textbf{the quality of test cases}
\item
  Number of test cases in BVA (normal) for 2 variables: 4n + 1 = \textbf{9}
\item
  DRE when 90 errors found before release and 10 after: \textbf{0.9}
\end{enumerate}

\textbf{Maintenance / SCM}

\begin{enumerate}
\def\labelenumi{\arabic{enumi}.}
\setcounter{enumi}{46}
\tightlist
\item
  Maintenance to accommodate a new OS: \textbf{adaptive}
\item
  Largest share of maintenance effort: \textbf{perfective}
\item
  Extracting design from source code: \textbf{reverse engineering}
\item
  A formally reviewed work product serving as basis for further development: \textbf{baseline}
\end{enumerate}

\textbf{More practice questions}

\begin{enumerate}
\def\labelenumi{\arabic{enumi}.}
\setcounter{enumi}{50}
\tightlist
\item
  In a DFD, a data store can connect directly to: \textbf{a process only}
\item
  Number of test cases in robust BVA for one variable: \textbf{7} (6n + 1 for n variables)
\item
  Worst-case BVA for 3 variables: 5³ = \textbf{125}
\item
  Intermediate COCOMO coefficient a for embedded mode: \textbf{2.8}
\item
  In UML, \texttt{\#} before an attribute means: \textbf{protected}
\item
  A hollow triangle arrowhead in a class diagram denotes: \textbf{generalisation (inheritance)}
\item
  In Musa\textquotesingle s basic model, failure intensity decreases: \textbf{linearly with the number of failures experienced}
\item
  Goel--Okumoto is an example of: \textbf{NHPP (non-homogeneous Poisson process) model}
\item
  Risk management process area is introduced at CMMI level: \textbf{3}
\item
  A decision table with 3 Boolean conditions has at most: \textbf{8 rules}
\item
  Which UML diagram best shows the order of messages between objects? \textbf{Sequence diagram}
\end{enumerate}

\begin{revbox}

\begin{itemize}
\tightlist
\item
  Unclear req → Prototype · Risk → Spiral · Changing req → Agile
\item
  Cohesion: Functional best, Coincidental worst · Coupling: Data best, Content worst
\item
  FP: VAF = 0.65 + 0.01ΣF (0.65--1.35), 14 GSCs · COCOMO organic 2.4/1.05/2.5/0.38
\item
  V(G) = E − N + 2 = decisions + 1 · BVA 4n+1
\item
  Alpha at dev site · Beta at customer site · Perfective maintenance largest
\end{itemize}

\end{revbox}
